
\noindent\emph{Proof.}
Assertion \textup{(i)} was already proved in \SGARefLink{sga3:xix:remark:2:4}{Exposé XIX, 2.4}. It
follows at once that \(N^\times\) is indeed a principal homogeneous
bundle under \(T\) for the laws defined in \textup{(ii)}; its local
triviality for the Zariski topology follows in particular from
\textup{\SGARefLink{sga3:xx:3.1:item:iv}{(iv)}}.

Let us prove \textup{(iii)}. If \(w\in N^\times(S)\), it is clear
that \(\alpha\circ\intaut(w)\) is a root of \(G\) relative to \(T\),
and hence is locally equal to \(\alpha\) or \(-\alpha\). Since it is
\(-\alpha\) on every fiber \textup{(Bible, 12-05, proof of the
corollary to Proposition 1)}, one has
\[
  \alpha\circ\intaut(w)=-\alpha.
\]
By transport of structure, it follows that
\[
  -\alpha^\ast
    =\intaut(w)^{-1}\circ\alpha^\ast
    =\intaut(w)\circ\alpha^\ast,
\]
since \(\intaut(w)^2=\intaut(w^2)\) and \(w^2\) is a section of
\(T\). Thus \(\intaut(w)\) induces inversion on \(T'\); since \(R\)
is central, it induces the identity on \(R\). The formula in
\textup{(iii)} defines a morphism \(T\to T\) having the same
properties, and therefore coincides with \(\intaut(w)\).

Let us prove \textup{\SGARefLink{sga3:xx:3.1:item:iv}{(iv)}}. Successively,
\[
\begin{aligned}
  w_\alpha(X)t w_\alpha(X)^{-1}
    &={}\exp(X)\exp(-X^{-1})\exp(X)t
       \exp(-X)\exp(X^{-1})\exp(-X)\\
    &={}\exp(X)\exp(-X^{-1})
       \exp\bigl(X-\alpha(t)X\bigr)
       \exp\bigl(\alpha(t)^{-1}X^{-1}\bigr)
       \exp\bigl(-\alpha(t)X\bigr)t.
\end{aligned}
\]
Applying formula \textup{(F)} gives
\[
\begin{aligned}
  \exp(-X^{-1})\exp\bigl((1-\alpha(t))X\bigr)
    ={}&\exp\bigl((\alpha(t)^{-1}-1)X\bigr)
       \alpha^\ast\bigl(\alpha(t)^{-1}\bigr)\\
     &{}\cdot\exp\bigl(-\alpha(t)^{-1}X^{-1}\bigr).
\end{aligned}
\]
Substitution in the preceding relation yields
\[
\begin{aligned}
  \intaut(w_\alpha(X))t
    &={}\exp\bigl(\alpha(t)^{-1}X\bigr)
       \alpha^\ast\bigl(\alpha(t)^{-1}\bigr)
       \exp\bigl(-\alpha(t)X\bigr)t\\
    &={}\exp(aX)\alpha^\ast\bigl(\alpha(t)^{-1}\bigr)t,
\end{aligned}
\]
where
\[
  a=\alpha(t)^{-1}
     -(\alpha\circ\alpha^\ast)\bigl(\alpha(t)^{-1}\bigr)\alpha(t).
\]
But \(\alpha\circ\alpha^\ast=2\), which gives \(a=0\) at once and
shows that \(w_\alpha(X)\in N^\times(S')\).

We now prove the second assertion of \textup{(iv)}. One has%
\[
\begin{aligned}
  \alpha^\ast(z)w_\alpha(X)
    &={}\exp(z^2X)\exp(-z^{-2}X^{-1})
       \exp(z^2X)\alpha^\ast(z)\\
    &={}\exp(zX)\exp\bigl((z^2-z)X\bigr)
       \exp(-z^{-2}X^{-1})\exp(z^2X)\alpha^\ast(z)\\
    &={}\exp(zX)\exp(-z^{-1}X^{-1})\alpha^\ast(z)^{-1}
       \exp\bigl((z^3-z^2)X\bigr)\exp(z^2X)\alpha^\ast(z)\\
    &={}\exp(zX)\exp(-z^{-1}X^{-1})\exp(zX)
       =w_\alpha(zX).
\end{aligned}
\]

Let us prove \textup{\SGARefLink{sga3:xx:3.1:item:v}{(v)}}. By the preceding result, the first
formula in \textup{\SGARefLink{sga3:xx:3.1:item:v}{(v)}} follows at once from the second; we prove
the latter:
\[
\begin{aligned}
  w_\alpha(X)^2
    &={}\exp(X)\exp(-X^{-1})\exp(2X)\exp(-X^{-1})\exp(X)\\
    &={}\exp(X)\exp(-X^{-1})\exp(X^{-1})
       \alpha^\ast(-1)\exp(-2X)\exp(X)\\
    &={}\exp(X)\alpha^\ast(-1)\exp(-X)
       =\alpha^\ast(-1),
\end{aligned}
\]
since
\[
  \alpha\bigl(\alpha^\ast(-1)\bigr)=(-1)^2=1.
\]
This also proves that
\(\alpha^\ast(-1)\in\operatorname{Centr}(G)(S)\).

Finally, let us prove \textup{\SGARefLink{sga3:xx:3.1:item:vi}{(vi)}}. The first assertion is a
special case of the second, so we prove the latter. Its two sides
define morphisms
\[
  W(\mathfrak g^\alpha)^\times
    \times_S W(\mathfrak g^{-\alpha})^\times
  \longrightarrow G.
\]
To prove that they coincide, it suffices to do so on a nonempty
open subset of every fiber \textup{(\SGARefLink{sga3:xviii:prop:1.4}{Exposé XVIII, 1.4})}. We may
therefore suppose that \(1+XY\) is invertible. Then, successively,
\[
\begin{aligned}
  w_\alpha(X)w_{-\alpha}(Y)
    ={}&\exp(X)\exp(-X^{-1})\exp(X)
        \exp(Y)\\
       &\qquad{}\exp(-Y^{-1})\exp(Y)\\
    ={}&\exp(X)\exp(-X^{-1})
        \exp\left(\frac{Y}{1+XY}\right)
        \alpha^\ast(1+XY)\\
       &\qquad{}
        \exp\left(\frac{X}{1+XY}\right)
        \exp(-Y^{-1})\exp(Y)\\
    ={}&\exp(X)
        \exp\left(\frac{-X^{-1}}{1+XY}\right)
        \alpha^\ast(1+XY)
        \exp\left(\frac{-Y^{-1}}{1+XY}\right)\exp(Y)\\
    ={}&\exp(-X^{-2}Y^{-1})
        \alpha^\ast\left(\frac{XY}{1+XY}\right)
        \exp(X+Y^{-1})\alpha^\ast(1+XY)\\
       &{}\qquad\cdot
        \exp\left(\frac{-Y^{-1}}{1+XY}\right)\exp(Y)\\
    ={}&\exp(-X^{-2}Y^{-1})\alpha^\ast(XY)
        \exp\left(\frac{Y^{-1}+X}{(1+XY)^2}\right)\\
       &\qquad{}
        \exp\left(\frac{-Y^{-1}}{1+XY}\right)\exp(Y)\\
    ={}&\alpha^\ast(XY)\exp(-Y)\exp(Y)
        =\alpha^\ast(XY).
\end{aligned}
\]

\paragraph{\textbf{Corollary 3.2.}}
\SGARefTarget{sga3:xx:3.2}
\emph{Let \(n\in\mathbb Z\), \(n\ne0\). For every \(w\in G(S)\),
the following conditions are equivalent:}

\smallskip
\noindent\emph{\textup{(i)} \(w\in N^\times(S)\);}

\smallskip
\noindent\emph{\textup{\SGAInlineRefTarget{sga3:xx:3.2:item:ii}(ii)}
\(\intaut(w)\circ n\alpha^\ast=-n\alpha^\ast\), where
\((n\alpha^\ast)(z)=\alpha^\ast(z)^n\).}

\smallskip
One has \textup{(i)}\(\Rightarrow\)\textup{\SGARefLink{sga3:xx:3.2:item:ii}{(ii)}} by \SGARefLink{sga3:xx:3.1}{Theorem 3.1}
\textup{\SGARefLink{sga3:xx:3.1:item:iii}{(iii)}}. Conversely, one may suppose that \(N^\times\) has a
section, and one is reduced to proving:

\paragraph{\textbf{Lemma 3.3.}}
\SGARefTarget{sga3:xx:3.3}
\emph{For \(n\ne0\), one has}
\[
  \operatorname{Centr}_G(n\alpha^\ast)=T.
\]

Indeed, the image \(T'\) of \(n\alpha^\ast\) is a subtorus of
\(G\). It follows from \SGARefLink{sga3:xix:paragraph:2:8}{Exposé XIX, 2.8} that
\(\operatorname{Centr}_G(n\alpha^\ast)\) is a reductive subgroup
of \(G\) containing \(T\). On every fiber,
\[
  \operatorname{Centr}_{G_s}(n\alpha^\ast_s)\ne G_s,
\]
and hence
\[
  \operatorname{Centr}_{G_s}(n\alpha^\ast_s)=T_s
\]
by \SGARefLink{sga3:xix:remark:1:6:3}{Exposé XIX, 1.6.3}.
Therefore
\[
  \operatorname{Centr}_G(n\alpha^\ast)=T,
\]
since these are smooth subgroups of \(G\).

\paragraph{\textbf{Remark 3.4.}}
\SGARefTarget{sga3:xx:3.4}
\emph{The construction of \(w_\alpha\), and the fact that
\(w_\alpha(X)\) normalizes \(T\), depend only on formula
\textup{(F)}. In particular, if \(G\) is an \(S\)-group satisfying
the conditions of \SGARefLink{sga3:xx:2.2}{2.2}, then \(\operatorname{Norm}_G(T)\) differs
from \(T\) on every fiber. Consequently, if \(G\) is an affine
\(S\)-group with connected fibers satisfying the conditions of \SGARefLink{sga3:xx:2.2}{2.2},
then \(G\) is reductive of semisimple rank \(1\). Indeed, it is
smooth in a neighborhood of the identity section, hence smooth, and
the criterion of \SGARefLink{sga3:xix:lemma:1:11}{Exposé XIX, 1.11} applies.}
